\section{Deletion-Risk Metric}

For a deletion target $z$ and an audit probe $q$, \method{} computes a normalized evidence vector:
\[
e(S_-,q,z) = [e_d,e_p,e_r,e_c,e_x,e_w],
\]
where $e_d$ measures direct leakage, $e_p$ paraphrased leakage, $e_r$ retrieval dependence, $e_c$ counterfactual answer dependence, $e_x$ extraction risk, and $e_w$ watermark or provenance-token hits. Each component lies in $[0,1]$.

\textbf{Operational scorers.}
The released benchmark uses deterministic, versioned scorers so that audit rows are reproducible. Direct leakage is normalized token overlap with the protected text. Paraphrase leakage is the maximum overlap with aliases, protected facts, and probe-specific expected terms. Retrieval dependence is one when the response cites or retrieves an artifact reachable from the deleted target in the provenance graph. Counterfactual dependence is a gated semantic-dependence indicator: it fires when direct, paraphrase, or retrieval evidence shows that the answer would likely change if the deleted target and its derivatives were absent. Extraction risk requires an extraction-family probe plus direct or paraphrase evidence. Watermark evidence is an exact canary or provenance-token hit. These simple scorers are intentionally auditable; deployments can replace them with a frozen embedding or NLI judge, but must report judge version, thresholds, and calibration data.

The per-target deletion risk is
\[
R_{\mathrm{del}}(z;S_-)=\mathbb{E}_{q\sim A(z)}[w^\top e(S_-,q,z)],
\]
where $A(z)$ is the target-specific audit-probe distribution and $w$ is a nonnegative weight vector. The default pilot weights are $0.25$ for direct leakage, $0.20$ for paraphrased leakage, $0.20$ for retrieval dependence, $0.15$ for counterfactual dependence, $0.15$ for extraction risk, and $0.05$ for watermark hits.

Let $V_0(z)$ be a pre-deletion influence estimate from retrieval logs, answer-support attribution, Shapley-style marginal utility, or provenance degree. \method{} defines a priority
\[
\pi(z)=\frac{\exp(\alpha V_0(z))}{\sum_{z'}\exp(\alpha V_0(z'))}.
\]
The aggregate request-level risk is
\[
R_{\method}(D_{\mathrm{del}};S_-)=\sum_{z\in D_{\mathrm{del}}}\pi(z)R_{\mathrm{del}}(z;S_-).
\]

Given sampled probes, \method{} reports an upper confidence bound using Hoeffding's inequality:
\[
\mathrm{UCB}(z)=\hat R_{\mathrm{del}}(z)+
\sqrt{\frac{\log(1/\delta)}{2n_z}}.
\]
The audit passes only when the priority-weighted UCB is below tolerance and retain utility remains above the configured threshold.

Because channel weights are policy choices, \method{} reports both the full evidence vector and the scalar risk. The default weights are not required for all applications: operators may choose semantic-heavy, retrieval-heavy, or extraction-heavy profiles, and the benchmark includes a weight/tolerance sensitivity table.

\textbf{Tolerance calibration.}
We recommend treating $\tau$ as an operating policy rather than a universal constant. A regulated or high-sensitivity domain should choose a low pass threshold and accept more escalations; an internal low-risk corpus may choose a higher threshold after validating false alarms on clean deletion cases. The artifact therefore reports detection, false alarms, utility, and audit calls across several tolerances instead of hiding calibration inside a single score.

\begin{proposition}[False-pass control]
Assume each probe evidence score $w^\top e(S_-,q,z)$ is independent and bounded in $[0,1]$, and assume the probe distribution $A(z)$ matches the declared audit threat model. If \method{} passes a target only when $\mathrm{UCB}(z)\leq \tau$, then the probability that the true target risk exceeds $\tau$ is at most $\delta$.
\end{proposition}

\textit{Proof sketch.}
Hoeffding's inequality gives
$\Pr[R_{\mathrm{del}}(z) > \hat R_{\mathrm{del}}(z)+\sqrt{\log(1/\delta)/(2n_z)}]\leq\delta$.
The pass condition requires the right-hand side to be no larger than $\tau$, so the probability of passing while $R_{\mathrm{del}}(z)>\tau$ is bounded by $\delta$ under the stated assumptions.

\begin{proposition}[Sample complexity]
Let a deletion target have true risk $R_{\mathrm{del}}(z)\geq\tau+\gamma$ for margin $\gamma>0$. If
\[
n_z \geq \frac{\log(1/\delta)}{2\gamma^2},
\]
then the UCB radius is at most $\gamma$, so an observed empirical risk near the true risk is sufficient to separate the target from the pass threshold.
\end{proposition}

These bounds are intentionally modest. They certify the audit process under a declared probe distribution; they do not prove absence of all possible future leakage. The main assumptions are bounded evidence scores, representative probe generation, stable post-deletion behavior during the audit, and access to enough response metadata to score the declared channels.

\begin{proposition}[Adaptive allocation control]
Let $X_t(z)\in[0,1]$ be the weighted evidence score observed at time $t$ for target $z$, and let probe choices be predictable with respect to the audit history. If the centered sequence $X_t(z)-\mathbb{E}[X_t(z)\mid \mathcal{F}_{t-1}]$ is a bounded martingale difference, then the same pass rule is valid with an anytime empirical-Bernstein or Hoeffding-Azuma radius in place of the i.i.d. Hoeffding radius.
\end{proposition}

This is the bound used to interpret valuation-guided allocation: adaptivity changes which target receives the next probe, but each scored response is still bounded and the selected probe is fixed before the response is observed. The artifact therefore reports the allocation policy and all row-level scores so that reviewers can recompute either the i.i.d. or martingale radius. Because probe families can be correlated, the reviewer-response artifact also reports seed-block bootstrap intervals, score-jitter stability, and the number of active evidence families per track.
